% GENERATED FILE. Edit canonical lessons, then run npm run generate:books.
% canonical-source-hash: fnv1a64:cb5dc2a1494abf9a
% canonical-lessons: MR-C555-avadhi, MR-C555-dirangai, MR-C555-tatdi, MR-C555-dhandal, MR-C555-tarambal

\chapter{Time Pressure}
\label{ch:mr-time-pressure}
\hlchaptermodality{555}

\begin{chapteropening}
\textbf{By the end of this chapter:} \emph{I can talk about time left, urgency and rushing.}

Complete the last lesson of chapter 555: I can talk about time left, urgency and rushing.
\end{chapteropening}

\section[\texorpdfstring{\mr{अवधी} (avadhī)}{अवधी (avadhī)}]{\mr{अवधी} (avadhī) — a span of time; time left}

\label{lesson:MR-C555-avadhi}



\begin{quote}
Before the new one: say the Marathi for mostly, then the Marathi for roughly.
\end{quote}

\subsection*{You'll want to know: \mr{अवधी}}
\textbf{\mr{अवधी}} — \emph{avadhī} — \textquotedblleft{}a span of time; time left\textquotedblright{}.

\begin{quote}
\mr{परीक्षेला} \mr{आता} \mr{फक्त} \mr{एक} \mr{आठवडा} \mr{अवधी} \mr{आहे}.
\end{quote}

\begin{quote}
\emph{parīkṣelā ātā phakta ek āṭhavḍā avadhī āhe.}
\end{quote}

\begin{quote}
There is only a week left before the exam.
\end{quote}

Masculine.

\begin{cousinweb}
\textbf{\mr{अवधी}} (\emph{avadhī}) is a tatsama, taken from Sanskrit \emph{avadhi}, \textquotedblleft{}a limit, a term\textquotedblright{} (Monier-Williams).
\end{cousinweb}

\begin{grammarlens}[title={What you've built}]
One more word for saying what happens if.
\end{grammarlens}

\subsection*{Your turn}
\begin{itemize}
\raggedright
  \item \emph{Say it:} \emph{avadhī}
  \item \emph{Say it:} \emph{avadhī}, once more
  \item \emph{Say it:} \emph{parīkṣelā ātā phakta ek āṭhavḍā avadhī āhe.} — There is only a week left before the exam.
  \item \emph{Recall:} say the Marathi for mostly, then the Marathi for roughly, then say \emph{avadhī} again
\end{itemize}

\begin{tcolorbox}[breakable,colback=teal!4,colframe=teal!35!black,title={Before you move on}]
What does \mr{अवधी} mean? (\textquotedblleft{}A span of time; time left\textquotedblright{}.) Say it once more.
\end{tcolorbox}
\section[\texorpdfstring{\mr{दिरंगाई} (diraṅgāī)}{दिरंगाई (diraṅgāī)}]{\mr{दिरंगाई} (diraṅgāī) — a needless delay; dilatoriness}

\label{lesson:MR-C555-dirangai}



\begin{quote}
Before the new one: say the Marathi for roughly, then the Marathi for a span of time.
\end{quote}

\subsection*{You'll want to know: \mr{दिरंगाई}}
\textbf{\mr{दिरंगाई}} — \emph{diraṅgāī} — \textquotedblleft{}a needless delay; dilatoriness\textquotedblright{}.

\begin{quote}
\mr{कामात} \mr{दिरंगाई} \mr{केली} \mr{तर} \mr{दंड} \mr{भरावा} \mr{लागतो}.
\end{quote}

\begin{quote}
\emph{kāmāt diraṅgāī kelī tar daṇḍ bharāvā lāgto.}
\end{quote}

\begin{quote}
If you put off the work, you have to pay a fine.
\end{quote}

Feminine.

\mr{दिरंगाई} \mr{करणे}, to drag one's feet; \mr{दिरंगाई} \mr{होणे}, a delay happens.

\begin{cousinweb}
\textbf{\mr{दिरंगाई}} (\emph{diraṅgāī}) is a loan, usually traced to Persian \emph{dirang}, \textquotedblleft{}delay, tarrying\textquotedblright{}, with the Marathi noun ending -\mr{आई}; the path into Marathi is not recorded in our sources.
\end{cousinweb}

\begin{grammarlens}[title={What you've built}]
One more word for saying what happens if.
\end{grammarlens}

\subsection*{Your turn}
\begin{itemize}
\raggedright
  \item \emph{Say it:} \emph{diraṅgāī}
  \item \emph{Say it:} \emph{diraṅgāī}, once more
  \item \emph{Say it:} \emph{kāmāt diraṅgāī kelī tar daṇḍ bharāvā lāgto.} — If you put off the work, you have to pay a fine.
  \item \emph{Recall:} say the Marathi for roughly, then the Marathi for a span of time, then say \emph{diraṅgāī} again
\end{itemize}

\begin{tcolorbox}[breakable,colback=teal!4,colframe=teal!35!black,title={Before you move on}]
What does \mr{दिरंगाई} mean? (\textquotedblleft{}A needless delay; dilatoriness\textquotedblright{}.) Say it once more.
\end{tcolorbox}
\section[\texorpdfstring{\mr{तातडी} (tātḍī)}{तातडी (tātḍī)}]{\mr{तातडी} (tātḍī) — urgency}

\label{lesson:MR-C555-tatdi}



\begin{quote}
Before the new one: say the Marathi for a span of time, then the Marathi for a needless delay.
\end{quote}

\subsection*{You'll want to know: \mr{तातडी}}
\textbf{\mr{तातडी}} — \emph{tātḍī} — \textquotedblleft{}urgency\textquotedblright{}.

\begin{quote}
\mr{कामाची} \mr{तातडी} \mr{होती}, \mr{म्हणून} \mr{तो} \mr{रविवारीही} \mr{कार्यालयात} \mr{गेला}.
\end{quote}

\begin{quote}
\emph{kāmācī tātḍī hotī, mhaṇūn to ravivārīhī kāryālayāt gelā.}
\end{quote}

\begin{quote}
The work was urgent, so he went to the office even on Sunday.
\end{quote}

Feminine.

\begin{cousinweb}
\textbf{\mr{तातडी}} (\emph{tātḍī}) is a tadbhava: CDIAL lists Marathi \emph{tā̃tḍī}, \textquotedblleft{}urgency, flurry\textquotedblright{}, under Sanskrit \emph{tānti}, \textquotedblleft{}suffocation\textquotedblright{} (CDIAL 5764).
\end{cousinweb}

\begin{grammarlens}[title={What you've built}]
One more word for saying what happens if.
\end{grammarlens}

\subsection*{Your turn}
\begin{itemize}
\raggedright
  \item \emph{Say it:} \emph{tātḍī}
  \item \emph{Say it:} \emph{tātḍī}, once more
  \item \emph{Say it:} \emph{kāmācī tātḍī hotī, mhaṇūn to ravivārīhī kāryālayāt gelā.} — The work was urgent, so he went to the office even on Sunday.
  \item \emph{Recall:} say the Marathi for a span of time, then the Marathi for a needless delay, then say \emph{tātḍī} again
\end{itemize}

\begin{tcolorbox}[breakable,colback=teal!4,colframe=teal!35!black,title={Before you move on}]
What does \mr{तातडी} mean? (\textquotedblleft{}Urgency\textquotedblright{}.) Say it once more.
\end{tcolorbox}
\section[\texorpdfstring{\mr{धांदल} (dhāndal)}{धांदल (dhāndal)}]{\mr{धांदल} (dhāndal) — a rush; a flurry}

\label{lesson:MR-C555-dhandal}



\begin{quote}
Before the new one: say the Marathi for a needless delay, then the Marathi for urgency.
\end{quote}

\subsection*{You'll want to know: \mr{धांदल}}
\textbf{\mr{धांदल}} — \emph{dhāndal} — \textquotedblleft{}a rush; a flurry\textquotedblright{}.

\begin{quote}
\mr{सकाळी} \mr{शाळेच्या} \mr{वेळी} \mr{घरात} \mr{खूप} \mr{धांदल} \mr{असते}.
\end{quote}

\begin{quote}
\emph{sakāḷī śāḷecyā veḷī gharāt khūp dhāndal aste.}
\end{quote}

\begin{quote}
In the morning at school time the house is in a great rush.
\end{quote}

Feminine.

\begin{cousinweb}
\textbf{\mr{धांदल}} (\emph{dhāndal}) is a Marathi word for a busy rush; its origin is not settled.
\end{cousinweb}

\begin{grammarlens}[title={What you've built}]
One more word for saying what happens if.
\end{grammarlens}

\subsection*{Your turn}
\begin{itemize}
\raggedright
  \item \emph{Say it:} \emph{dhāndal}
  \item \emph{Say it:} \emph{dhāndal}, once more
  \item \emph{Say it:} \emph{sakāḷī śāḷecyā veḷī gharāt khūp dhāndal aste.} — In the morning at school time the house is in a great rush.
  \item \emph{Recall:} say the Marathi for a needless delay, then the Marathi for urgency, then say \emph{dhāndal} again
\end{itemize}

\begin{tcolorbox}[breakable,colback=teal!4,colframe=teal!35!black,title={Before you move on}]
What does \mr{धांदल} mean? (\textquotedblleft{}A rush; a flurry\textquotedblright{}.) Say it once more.
\end{tcolorbox}
\section[\texorpdfstring{\mr{तारांबळ} (tārāmbaḷ)}{तारांबळ (tārāmbaḷ)}]{\mr{तारांबळ} (tārāmbaḷ) — a scramble; a frantic hurry}

\label{lesson:MR-C555-tarambal}



\begin{quote}
Before the new one: say the Marathi for urgency, then the Marathi for a rush.
\end{quote}

\subsection*{You'll want to know: \mr{तारांबळ}}
\textbf{\mr{तारांबळ}} — \emph{tārāmbaḷ} — \textquotedblleft{}a scramble; a frantic hurry\textquotedblright{}.

\begin{quote}
\mr{पाहुणे} \mr{लवकर} \mr{आले}, \mr{म्हणून} \mr{आमची} \mr{तारांबळ} \mr{उडाली}.
\end{quote}

\begin{quote}
\emph{pāhuṇe lavkar āle, mhaṇūn āmcī tārāmbaḷ uḍālī.}
\end{quote}

\begin{quote}
The guests came early, so we were thrown into a scramble.
\end{quote}

Feminine.

\mr{तारांबळ} \mr{उडणे}, to be thrown into a scramble.

\begin{cousinweb}
\textbf{\mr{तारांबळ}} (\emph{tārāmbaḷ}) is a Marathi word for a frantic scramble; its origin is not settled.
\end{cousinweb}

\begin{grammarlens}[title={What you've built}]
One more word for saying what happens if.
\end{grammarlens}

\subsection*{Your turn}
\begin{itemize}
\raggedright
  \item \emph{Say it:} \emph{tārāmbaḷ}
  \item \emph{Say it:} \emph{tārāmbaḷ}, once more
  \item \emph{Say it:} \emph{pāhuṇe lavkar āle, mhaṇūn āmcī tārāmbaḷ uḍālī.} — The guests came early, so we were thrown into a scramble.
  \item \emph{Recall:} say the Marathi for urgency, then the Marathi for a rush, then say \emph{tārāmbaḷ} again
\end{itemize}

\begin{tcolorbox}[breakable,colback=teal!4,colframe=teal!35!black,title={Before you move on}]
What does \mr{तारांबळ} mean? (\textquotedblleft{}A scramble; a frantic hurry\textquotedblright{}.) Say it once more.
\end{tcolorbox}
